\documentclass{amsart}
\usepackage[utf8x]{inputenc}
\usepackage[english]{babel}


\usepackage{titlesec}
\titleformat{\section}{\normalfont\bf\Large}{\relax }{0em}{}
\titleformat{\subsection}{\normalfont\bf\large }{\relax}{0em}{}
\titleformat{\subsubsection}{\normalfont\bf}{\relax}{0em}{}

\usepackage{enumitem}
\usepackage{quoting}
\usepackage{array}
\usepackage{xcolor}
\long\def\gray#1{{\color{gray}{#1}}}\let\grey=\gray
\long\def\blue#1{{\color{blue}{#1}}}
\long\def\red#1{{\color{red}{#1}}}
\long\def\green#1{{\color{green}{#1}}}
\long\def\violet#1{{\color{violet}{#1}}}
\long\def\lila#1{{\color{lila}{#1}}}
\long\def\cyan#1{{\color{cyan}{#1}}}

\usepackage{wrapfig}
\usepackage{lpic}

\usepackage{manyfoot}
\SetFootnoteHook{\hspace{1.5em}}
\DeclareNewFootnote{A}
\SetFootnoteHook{\hspace{1.5em}}
\DeclareNewFootnote{E}[alph]
\usepackage[hyperfootnotes=false,
    colorlinks,
    linkcolor={blue!80!black},
    citecolor={blue!50!black},
    urlcolor={blue!80!black}]{hyperref}
\def\foothref#1#2{\href{#1}{ #2}\footnoteA{\url{#1}}}

\newcommand{\skipthis}[2][\dots]{\red{[#1]}\gray{#2}}
\long\def\skipthis#1{\relax}

\newtheorem*{theorem}{Theorem}
\newtheorem*{proposition}{Proposition}
\newtheorem*{question}{Question}
\newtheorem*{conjecture}{Conjecture}
\newtheorem*{definition}{Definition}

\let\dmo=\DeclareMathOperator

\let\tmp\epsilon
\let\epsilon\varepsilon
\let\varepsilon\tmp


\begin{document}
\centerline{\bf\sc\Large FROM TOPOGRAPHY TO GEOMETRY%
  \footnoteE{English translation by Rostislav Matveev of the article:
    Patrick Popescu-Pampu,  \textit{De la topographie \`a la
      g\'eom\'trie I et II}, Images des
    Math\'ematiques (2017);\\
    \url{https://doi.org/10.60868/7xv4-9g98} and
    \url{https://doi.org/10.60868/tzg9-3c39} (Ed.)}}%
\smallskip\par
\centerline{\sc\large Contribution of Maps to Mathematical Thinking}
\bigskip\par

\centerline{by Patrick Popescu-Pampu}

\bigskip\par

\begin{center}
  \sc\large
  PART ONE:\\
  From topographic maps to the structure of surfaces
\end{center}
\bigskip\par\noindent
Did you know that Paris is a city of eight hills? Here is a map that
lets you see it:

\begin{center}
  \begin{lpic}[draft,b(6ex),t(2ex),clean]{Pix/Paris-elevation(0.65)}
     \lbl[t]{73,-5;\parbox{0.72\textwidth}{
         \begin{center}\sf
           The city of eight hills.
         \end{center}}}
  \end{lpic}
\end{center}
On the map above one can see the names of these hills, as well as
their elevations. But one also observes a color gradient, with the
darker shades corresponding to the higher levels. The boundary between
two regions of different colors is a line of constant elevation. A map
on which such lines are shown is called ``topographic'' --- from the
Greek
% τόπος
$\tau o\pi o\zeta$
(place) and
%γραφία
$\gamma\rho\alpha\phi\iota\alpha$ (to draw).

Here is another one, in which the collection of constant-elevation
lines is supplemented with information about roads, buildings, and
municipalities:
\begin{center}
  \begin{lpic}[draft,b(4ex),t(1ex),clean]{Pix/Fg-Mons(0.25)}
     \lbl[t]{180,-10;\parbox{0.75\textwidth}{
         \begin{center}\sf
           A topographic map supplemented with additional information.
         \end{center}}}
  \end{lpic}
\end{center}

How does one construct a map of this type? The theoretical principle
is to represent the landscape first by a smooth surface, obtained by
smoothing out overly pronounced indentations --- much as if a heavy
snowfall had covered it. One then slices this surface by horizontal
planes at different heights, taken successively at equal distances
from one another, and then projects the resulting intersection curves
onto a single horizontal plane. This process is illustrated in the
following figure:

\begin{center}
  \begin{lpic}[draft,b(8ex),t(1ex),clean]{Pix/Construction-nbg(0.35)}
     \lbl[t]{90,-10;\parbox{0.75\textwidth}{
         \begin{center}\sf
           Principle for constructing a topographic map, illustrated
           using a smooth surface representing a landscape with two
           hills.
         \end{center}}}
  \end{lpic}
\end{center}
Of course, since the Earth is not flat, the notion of a ``horizontal''
plane is not well defined. In any case, the idea is that these
intersection curves are close to lines of constant elevation, which is
why their projections are called ``contour lines.'' Note that the
intersection of the smooth surface representing the landscape with a
horizontal plane generally consists of several separate pieces, or, as
mathematicians call them, ``connected components.'' Here, a ``contour
line'' refers to one of these connected components.

In its simplest form, which does not use colors, a topographic map is
such a collection of contour lines. Most of them are smooth,%
\footnoteA{Intuitively, a ``smooth'' curve has, at each of its points,
  a well-defined tangent line. In current mathematical usage, the term
  ``smooth'' is in fact even more restrictive: one additionally
  requires that the slopes of the tangents form a function, which is
  itself smooth, that is whose graph is a new curve that itself has a
  tangent at each point, and that these tangents again have the same
  property, and so on ad infinitum. These subtleties will not matter
  here.} %
but some are reduced to points --- summits and pits --- and others
intersect themselves.%
\footnoteE{Still even more complicated behaviour is also possible ---
  for example if the landscape contains a completely flat valley and
  one of the slicing planes lies exactly at the valley's height. In
  this case the ``counter line'' is not a line at all, but includes a
  whole region of the map. (Ed.)} %
In the previous figure, such self-intersection happens for the one
obtained by taking the plane that passes through the saddle, that is,
the second horizontal plane counting from the bottom.

The aim of this two-part article is to explain that mathematical
thinking about the structure of topographic maps has been remarkably
fertile in important discoveries about the qualitative --- so-called
topological --- shape of “spaces” of any dimension. We shall see, among
other things, that this line of thought provides one of the main tools
for exploring the topological structure of such spaces. I also intend
to convey some of the reasons that led mathematicians to go beyond the
three dimensions of everyday intuition, in order to explore
qualitatively ``spaces'' even of infinite dimension.

\section{A few reflections by Cayley}
The story begins in 1859, when the British mathematician Arthur Cayley
wrote an article%
\footnoteA{ Arthur Cayley, \textit{On contour and slope lines,}
  Philosophical Magazine 18 (1859), pp. 264–268.} %
in which he wondered about the possible shapes of contour lines, and
about their qualitative change as one varies the elevation:
\begin{quotation}\sf
  CITE
\end{quotation}

For example, in the preceding drawing there are two ``summits'' and
one ``knot,'' but no ``immet.'' From now on, I will write ``bottom''
rather than ``immet.''

``Summits,'' ``knots,'' and ``bottoms'' are called the ``singular
points'' of the topographic map. Near the other points, which
mathematicians usually call ``generic points'' of the map, the contour
lines look like a family of slightly distorted parallel lines:

\begin{center}
  \begin{lpic}[draft,b(8ex),t(1ex),clean]{Pix/Voisinage-nbg-en(0.45)}
     \lbl[t]{90,-10;\parbox{0.75\textwidth}{
         \begin{center}\sf
          XXX \red{REDO}
         \end{center}}}
  \end{lpic}
\end{center}




















































\end{document}
